% Source: OpenAI2026AreaLaw, 09-amplification.tex, lines 547--561, and
% 01-preliminaries.tex, lines 33--46 and 76--80, revision
% openai/math@adc7f1241b42e322a6451854ab7e4b4c146bf78a.

\section{The final entropy comparison of the amplification}

The proof of \cite[Proposition~10.2]{OpenAI2026AreaLaw} ends with a unitary
$U'$ that acts on neither the target system nor a distant region $J$ and that
moves the state close to a vector whose joint reduction on these two systems
is a product. The results of this section turn that unitary estimate into the
bound $C n^{-M}$ on mutual information. They form the last step of the proof
and do not by themselves establish the proposition: the construction of the
unitary is a separate argument.

Throughout, $V$ is a finite set of sites with finite local spaces
$\mathcal H_v$, and $\mathcal H=\bigotimes_{v\in V}\mathcal H_v$. For a region
$R\subseteq V$ put $\mathcal H_R=\bigotimes_{v\in R}\mathcal H_v$, and for a
vector $\psi\in\mathcal H$ let $\rho_{\psi,R}$ be the reduced state of
$|\psi\rangle\langle\psi|$ on $R$. For disjoint regions $X,J$ write
$I_\psi(X:J)=S(\rho_{\psi,X})+S(\rho_{\psi,J})-S(\rho_{\psi,X\cup J})$.
Ancillary systems are included by enlarging $V$.

\begin{lemma}[Trace distance of pure states]
    \label{lem:al_amp_pure_trace_distance}
    \lean{TNLean.PEPS.AreaLaw.traceDistance_vecMulVec_le_norm_sub,
        FiniteProduct.traceDistance_reducedPure_le}
    \leanok
    For unit vectors $\psi,\varphi$ of a finite-dimensional space,
    \begin{align*}
        \tfrac12\bigl\||\psi\rangle\langle\psi|
        -|\varphi\rangle\langle\varphi|\bigr\|_1 & \le\|\psi-\varphi\|.
    \end{align*}
    For unit vectors of $\mathcal H$ the reduced states on every region
    $R\subseteq V$ are no farther apart in trace distance than the two pure
    states.
\end{lemma}
\begin{proof}\leanok
    The trace distance of the two pure states is
    $\sqrt{1-|\langle\psi,\varphi\rangle|^2}$. With
    $t=|\langle\psi,\varphi\rangle|\le1$ we have
    $1-t^2\le2-2t\le2-2\Re\langle\psi,\varphi\rangle
        =\|\psi-\varphi\|^2$.
    Partial trace contracts trace distance.
\end{proof}

\begin{lemma}[A unitary away from a region]
    \label{lem:al_amp_unitary_complement}
    \lean{FiniteProduct.reducedPure_unitary_compl,
        FiniteProduct.mutualInformation_eq_of_reducedPure_union_eq,
        FiniteProduct.trace_reducedPure_eq_norm_sq}
    \leanok
    Let $R\subseteq V$ and let $U$ be a unitary on $\mathcal H_{V\setminus R}$.
    For every $\varphi\in\mathcal H$, the vectors $\varphi$ and
    $(I_R\otimes U)\varphi$ have the same reduced state on $R$. In particular,
    for disjoint regions $X,J$ with $R=X\cup J$, they have the same mutual
    information $I(X:J)$.
\end{lemma}
\begin{proof}\leanok
    The reduced state of $(I_R\otimes U)\varphi$ is the partial trace over
    $V\setminus R$ of $(I\otimes U)|\varphi\rangle\langle\varphi|(I\otimes
        U)^\dagger$. Cyclicity of the partial trace in the traced factor and
    $U^\dagger U=I$ remove $U$. The mutual information of $X$ and $J$ is a
    function of the reduced state on $X\cup J$.
\end{proof}

\begin{lemma}[Continuity against a product reduction]
    \label{lem:al_amp_continuity_product}
    \lean{FiniteProduct.mutualInformation_le_of_norm_sub_le}
    \leanok
    \uses{lem:al_amp_pure_trace_distance}
    Let $\psi,\eta\in\mathcal H$ be unit vectors and let $X,J\subseteq V$ be
    disjoint with $I_\eta(X:J)=0$. Then
    \begin{align*}
        I_\psi(X:J) & \le 8\sqrt{\|\psi-\eta\|}\,\log(e\dim\mathcal H_X).
    \end{align*}
    Only the dimension of $\mathcal H_X$ enters.
\end{lemma}
\begin{proof}\leanok
    Apply \cite[Lemma~2.1]{OpenAI2026AreaLaw} to $S(X)$ and $S(X|J)$ for the
    joint reduced states on $X\cup J$, and use $I(X:J)=S(X)-S(X|J)$. The two
    moduli sum to $8\sqrt\delta\log(e\dim\mathcal H_X)$, where $\delta$ is the
    trace distance of the joint reductions. By
    Lemma~\ref{lem:al_amp_pure_trace_distance}, $\delta\le\|\psi-\eta\|$.
\end{proof}

\begin{theorem}[Entropy comparison after the local unitary]
    \label{thm:al_amp_entropy_comparison}
    \lean{FiniteProduct.mutualInformation_le_of_unitary_sub_le}
    \leanok
    \uses{lem:al_amp_unitary_complement, lem:al_amp_continuity_product}
    Let $X,J\subseteq V$ be disjoint, let $U$ be a unitary on
    $\mathcal H_{V\setminus(X\cup J)}$, and let $\varphi,\eta\in\mathcal H$
    be unit vectors with $I_\eta(X:J)=0$. Put
    $\delta=\|(I_{X\cup J}\otimes U)\varphi-\eta\|$. Then
    \begin{align*}
        I_\varphi(X:J) & \le 8\sqrt\delta\,\log(e\dim\mathcal H_X).
    \end{align*}
    This is the comparison in the final part of the proof of
    \cite[Proposition~10.2]{OpenAI2026AreaLaw}.
\end{theorem}
\begin{proof}\leanok
    The vector $(I\otimes U)\varphi$ is a unit vector with the same mutual
    information $I(X:J)$ as $\varphi$, by
    Lemma~\ref{lem:al_amp_unitary_complement}. Apply
    Lemma~\ref{lem:al_amp_continuity_product} to it and $\eta$.
\end{proof}

\begin{theorem}[Polynomial decay of the amplified information]
    \label{thm:al_amp_entropy_comparison_scale}
    \lean{FiniteProduct.mutualInformation_le_rpow_of_unitary_sub_le,
        FiniteProduct.mutualInformation_le_inv_pow_of_unitary_sub_le}
    \leanok
    \uses{thm:al_amp_entropy_comparison}
    In the setting of Theorem~\ref{thm:al_amp_entropy_comparison}, let
    $n\ge1$ be an integer and suppose that $\delta\le C_1n^{-K/2}$ and
    $\log\dim\mathcal H_X\le C_2n^2$. Then
    \begin{align*}
        I_\varphi(X:J) & \le 8\sqrt{C_1}\,(1+C_2)\,n^{2-K/4}.
    \end{align*}
    If $M,K$ are integers with $K\ge4(M+2)$, then every $X'\subseteq X$
    satisfies $I_\varphi(X':J)\le8\sqrt{C_1}\,(1+C_2)\,n^{-M}$. The constant
    depends only on $C_1$ and $C_2$. The choice $K\ge4(M+3)$ in the proof of
    \cite[Proposition~10.2]{OpenAI2026AreaLaw} satisfies this condition.
\end{theorem}
\begin{proof}\leanok
    The square root of $C_1n^{-K/2}$ is $\sqrt{C_1}\,n^{-K/4}$, and
    $\log(e\dim\mathcal H_X)\le1+C_2n^2\le(1+C_2)n^2$ since $n\ge1$.
    Multiply the two bounds. Since $\log\dim\mathcal H_X\ge0$, the constant
    $C_2$ is nonnegative, and $2-K/4\le-M$ gives $n^{2-K/4}\le n^{-M}$.
    Discarding part of $X$ does not increase mutual information.
\end{proof}
